
\subsubsection{Sub-Quadratic Sequence Models}

The quadratic complexity of transformer attention has motivated efficient alternatives. Longformer and BigBird introduce sparse attention patterns reducing complexity to O(n). Linear attention methods reformulate attention as kernel feature maps. More recently, Mamba introduces selective state spaces with input-dependent gating, achieving O(n) complexity. RWKV combines linear attention with RNN-style recurrence. These architectures achieve efficiency but provide no mechanism for task-specific adaptation during deployment.

\subsubsection{Efficient Adaptation Methods}

LoRA learns low-rank updates to pretrained weights. Adapters insert bottleneck modules between layers. Prompt tuning optimizes soft prompts. These methods assume the base model preserves adaptation capability, an assumption that may not hold after architecture conversion. LoRA applied post-conversion must work with potentially degraded representations.

\subsubsection{Model Conversion and Distillation}

Knowledge distillation transfers knowledge by matching output distributions. For architecture conversion, distillation typically optimizes KL divergence on logits and MSE on hidden states. StreamingLLM and H2O focus on KV cache compression for efficient inference. The conversion literature optimizes for output fidelity, not adaptation preservation. A converted model may match teacher accuracy on fixed tasks while losing the ability to adapt to new ones.

\subsubsection{Positioning}

Prior work addresses efficiency (Mamba, RWKV), adaptation (LoRA, adapters), or output preservation (distillation) in isolation. TC-SSM integrates task conditioning into the conversion process itself, aiming to enable sub-quadratic models that retain few-shot capability.

